
@@M0@@
où @@M1@@ est obtenu à partir de @@M2@@ par changement de base ; ce changement de base étant plat, @@M3@@ est une immersion régulière, et par suite, @@M4@@ est une immersion régulière de @@M5@@ dans un schéma lisse sur @@M6@@, d'où le résultat.

\textbf{Proposition 1.6.} \uline{Soit} @@M7@@ \uline{un morphisme et} @@M8@@ \uline{un changement de base plat}. \uline{Si} @@M9@@ \uline{est d'intersection complète}, \uline{alors le morphisme} @@M10@@ \uline{déduit de} @@M11@@ \uline{l'est aussi}. \uline{La réciproque est vraie si} @@M12@@ \uline{est de plus quasi-compact et surjectif}.

On peut trouver un recouvrement de @@M13@@ par des ouverts @@M14@@ tels que pour tout @@M15@@, la restriction de @@M16@@ à @@M17@@ se factorise en @@M18@@, où @@M19@@ est une immersion régulière, et @@M20@@ un morphisme lisse. Les @@M21@@ forment un recouvrement ouvert de @@M22@@, les immersions @@M23@@ sont régulières d'après VII 1.5, et les morphismes @@M24@@ sont lisses, d'où le premier résultat. Pour la réciproque on conclut d'abord de EGA IV 2.7.1 (iv) que @@M25@@ est localement de présentation finie, ce qui nous ramène facilement, grâce à 1.2, au cas où @@M26@@ est une immersion fermée, auquel cas l'assertion n'est autre que VII 1.5.

\textbf{Proposition 1.7.} \uline{Soit} @@M27@@ \uline{un morphisme d'intersection complète}. \uline{Alors} @@M28@@ \uline{est un morphisme parfait}.

La propriété étant locale sur @@M29@@, on peut se placer au voisinage d'un point @@M30@@ de @@M31@@, et supposer qu'il existe un ouvert @@M32@@ contenant @@M33@@, et une factorisation :
